\chapter{Disjunction Makes a Choice}
\label{ch:disjunction}

A constructive proof of $A\lor B$ does more than certify that one side holds. It exposes which side was constructed and carries the evidence needed by that branch. In $\mathrm{CH}\text{-}0$, this is a sum type.

\begin{formal}[title={Sum rules}]
\[
\frac{\Gamma\vdash a:A}{\Gamma\vdash\operatorname{inl}(a):A+B}\;+I_1
\qquad
\frac{\Gamma\vdash b:B}{\Gamma\vdash\operatorname{inr}(b):A+B}\;+I_2
\]
\[
\frac{\Gamma\vdash s:A+B\quad\Gamma,x{:}A\vdash t:C\quad\Gamma,y{:}B\vdash u:C}
{\Gamma\vdash\operatorname{case}\;s\;\operatorname{of}\;\operatorname{inl}(x)\Rightarrow t\mid\operatorname{inr}(y)\Rightarrow u:C}\;+E.
\]
\end{formal}

Case analysis is typed branching: the scrutinee chooses a branch, but both branches must converge on the same result type $C$.

\begin{figure}[p]\centering\begin{tikzpicture}[gclplate]\node[anchor=west,font=\sffamily\bfseries\Large,gcllabel] at (-6,3.8){DISJUNCTION MAKES A CHOICE};\draw[GCLWarm,line width=1pt](-6,3.45)--(6,3.45);\node[gcllabel,draw](s) at (0,2){$A+B$};\node[gcllabel,draw](l) at (-3,.2){$\operatorname{inl}(a)$};\node[gcllabel,draw](r) at (3,.2){$\operatorname{inr}(b)$};\draw[gclwarmarrow](s)--(l);\draw[gclwarmarrow](s)--(r);\end{tikzpicture}\caption{\textbf{Plate 10. Disjunction makes a choice.} A constructive sum inhabitant records a branch and its payload. Scope: richer than a proof-irrelevant Boolean, but not a universal semantics of disjunction.}\label{plate:disjunction-choice}\end{figure}


\begin{lemma}[Sum canonical form]
If $v$ is closed, $\vdash v:A+B$, and $v$ is a value, then either $v=\operatorname{inl}(v_A)$ with $\vdash v_A:A$, or $v=\operatorname{inr}(v_B)$ with $\vdash v_B:B$.
\end{lemma}
\begin{proof}
By inspection of the value grammar and inversion of typing. No other value constructor concludes with a sum type.
\end{proof}

\begin{theorem}[Principal case preservation]
A well-typed case expression whose scrutinee is a matching injection contracts to a well-typed selected branch of the same result type.
\end{theorem}
\begin{proof}
For the left case, inversion gives $\Gamma\vdash a:A$ and $\Gamma,x{:}A\vdash t:C$. The reduct is $t[a/x]$, typed $C$ by substitution. The right case is symmetric.
\end{proof}

\begin{figure}[p]\centering\begin{tikzpicture}[gclplate]\node[anchor=west,font=\sffamily\bfseries\Large,gcllabel] at (-6,3.8){CASE BRANCHES CONVERGE};\draw[GCLWarm,line width=1pt](-6,3.45)--(6,3.45);\node[gcllabel,draw](l) at (-3,1.7){$x:A\vdash t:C$};\node[gcllabel,draw](r) at (3,1.7){$y:B\vdash u:C$};\node[gcllabel,draw](c) at (0,.1){case result $:C$};\draw[gclbluearrow](l)--(c);\draw[gclbluearrow](r)--(c);\end{tikzpicture}\caption{\textbf{Plate 11. Case branches converge.} Sum elimination may inspect either constructor, but both branches must produce the same result type. Scope: the stated $\mathrm{CH}\text{-}0$ rule.}\label{plate:case-converge}\end{figure}


\begin{workedexample}[title={Choice is evidence}]
The term
\[
\lambda x{:}P.\operatorname{inl}(x):P\to(P+Q)
\]
not only proves the disjunction; it records that the left alternative was chosen. A consumer can inspect that constructor and run the corresponding branch. This computational choice is precisely what a constructive disjunction proof contributes.
\end{workedexample}

\begin{figure}[p]\centering\begin{tikzpicture}[gclplate]\node[anchor=west,font=\sffamily\bfseries\Large,gcllabel] at (-6,3.8){SUM DETOUR};\draw[GCLWarm,line width=1pt](-6,3.45)--(6,3.45);\node[gcllabel,draw](r) at (0,2){case $(\operatorname{inl}a)$};\node[gcllabel,draw](d) at (0,.2){selected branch $t[a/x]$};\draw[gclwarmarrow](r)--(d);\end{tikzpicture}\caption{\textbf{Plate 12. Sum detour.} Injection followed immediately by case analysis contracts to the selected branch by substitution. Scope: principal reduction only; no arbitrary-extension confluence claim.}\label{plate:sum-detour}\end{figure}


\begin{warningbox}
A sum inhabitant is not merely a Boolean bit. It carries a tag and a payload. Conversely, this chapter does not assert that every logical disjunction in every semantics should be implemented as a tagged union.
\end{warningbox}

\begin{labbox}
Run \texttt{python labs/lab04_sums.py}. It evaluates retained left/right case fixtures, verifies branch-result type agreement, and rejects mismatched branch codomains.
\end{labbox}

\section*{Exercises}\addcontentsline{toc}{section}{Exercises}
\begin{enumerate}[label=\textbf{4.\arabic*.}]
\item \textbf{Checkpoint.} State both sum introduction rules.
\item \textbf{Checkpoint.} Why must both case branches have type $C$?
\item \textbf{Checkpoint.} Contract a case over $\operatorname{inr}(b)$.
\item \textbf{Core.} Derive $P\to(P+Q)$.
\item \textbf{Core.} Derive $(P+Q)\to(Q+P)$.
\item \textbf{Core.} Type a case expression that maps $P+Q$ into a common type $R$ given $P\to R$ and $Q\to R$.
\item \textbf{Synthesis.} Contrast sum evidence with a Boolean truth flag.
\item \textbf{Synthesis.} Explain how case analysis exposes proof relevance.
\item \textbf{Proof Workshop.} Prove substitution for the case constructor, including both binders.
\item \textbf{Proof Workshop.} Prove principal case preservation for the right injection.
\item \textbf{Design Clinic.} Specify an AST and type-checking rule for case expressions.
\item \textbf{Challenge.} State a sum eta principle and classify it as operational, equational, or semantic under the conventions of this volume.
\end{enumerate}
